\section{Strip crossing}\label{L:part6}

We calculate the total critical mass of bridges across a honeycomb strip.
Number the parallel triangle-side grid lines in any fixed lattice direction
at unit spacing. A strict bridge of height $H\ge1$ joins a fixed midport on
level $0$ to a midport on level $H$, with every visited triangle center
strictly between the two lines. Its terminal port is otherwise unrestricted.
Write $B_H$ for the sum of $X^{|\omega|}$ over these bridges, where
$|\omega|$ counts visited centers, and put $B_0=1$.

\begin{proposition}[Open-strip crossing mass]\label{L:prop:strip-mass}
There are constants $0<c<C<\infty$ such that, for every integer $H\ge1$,
\[
 cH^{-1/4}\le B_H\le CH^{-1/4}.
\]
The bound holds in each of the three lattice-side directions.
\end{proposition}

The proof expresses lower half-strip states by polynomials and extracts
a determinant for the crossing mass. Positive column comparison relates
the homogeneous square-rhombus point to physical honeycomb width. We then
evaluate the determinant at that square point by a positive Gram matrix
and a three-column Hermite remainder. Positivity and strip exhaustion
identify the finite determinant identities with the infinite-strip mass.

Keep the cells, ports, weights, swaps \(B(s)\), triangle/cap maps \(P,C\) and loop-killing diagram calculus of the periodic finite calculation above, but do not identify the two side boundaries. Use \(n\) columns extending along all integer rows. In this section recenter the rapidities: column \(j\) has row-cell argument \(s=p z_j,\ p=q^{3/2}=e^{3\pi i/8}\); spatial swaps still use the right \(z\) divided by the left. Put
\[
\begin{gathered}
 [z]=(z+z^{-1})/2,\quad x_j=[z_j],\\
 a=\sin(\pi/8),\quad b=\cos(\pi/8),\quad s_0=a+b,\quad d=b-a,\\
 \tau=2a,\quad \gamma=1-\tau^2=\sqrt2-1 .
\end{gathered}
\]
Thus \(d=\sqrt2 a,\ s_0=\sqrt2 b,\ ab=\sqrt2/4\). Honeycomb (\(g=\pi/3\) in the cell parametrization) is \(z=q^{-1/2},x=b\), and square rhombi give \(z=x=1\). Vertical reflection changes \(z\) to \(z^{-1}\), swapping \(u,U\) and \(w,W\). Useful weight identities (with coherent half-powers) are
\[
\begin{gathered}
 v=\frac{x-a}{x+b},\quad d_z=1-v=\frac{s_0}{x+b},\\
 g_z=u/d_z,\quad h_z=U/d_z=g_{1/z},\\
 d_z g_z h_z=\tau v,\quad g_z^2+h_z^2+\tau g_z h_z=1,\\
 w=d_z g_z^2-\gamma v,\quad W=d_z h_z^2-\gamma v .
\end{gathered}
\]
Indeed with \(\eta\) as in the cell definition, the sine denominator is \(-(x+b)/2\), and \(g_z=\sin(5\pi/8+\eta)/b,\ h_z=\sin\eta/b\), linear combinations of \(z^{\pm1/2}\); the identities follow by the addition formulas.

At a slice (e.g. just over row 0), let \(E\) be the vector of lower half-strip arc sums with no side contacts and ends on the slice, and \(O\) the same with exactly one end on the left side (summed over rows up to 0). No loops. States are noncrossing pairings of subsets of the ordered ports, with in the odd case a single occupied unpaired port attached to the left side, called the defect (no pair spans it). Row addition with sides vacant is denoted \(T\), on either block; adding a row with a left source only, from even to odd, is \(A_L\). Thus
\[
 E_\emptyset=1,\quad E=T E,\qquad O=T O+A_L E .
\]

These sums converge, and the systems with this normalization are nonsingular, at either homogeneous honeycomb angle (\(s=q\) or \(q^2\)) and locally around each. Indeed by the boundary flux identity proved above on convex polygons of equilateral triangles, each boundary source of a finite parallelogram of these cells has bounded total positive \(X\)-weight of simple exit walks (first exit, at mid-edges), uniformly in the dimensions: boundary coefficients are \(\cos(3\,{\rm turn}/8)\) with \(|{\rm turn}|\le\pi\). Thus on exhaustion individual arc sums in the lower half-strip are bounded (also from a fixed top port to a variable left port), and one may drop mutual avoidance. Every allowed state is accessible from empty even by rows, using once a source row for odd: working downwards, join a pair consecutive among occupied sites in a fresh row using corners and horizontals, continuing others vertically, repeat and finally route the single defect to the left. Single arc weights are strictly positive. Since the empty output on the source-free even block sees only empty input, convergence of the positive series for each nonempty even and each odd state and this accessibility imply summability of powers on the nonempty even block and on the odd block. Both spectral radii are \(<1\). This yields also local absolute convergence and rational continuation (in chosen square roots). Braid intertwining with \(T,A_L\) and inversion, using row addition from left to right, imply the rational covariance of the actual \(E,O\): exchanging adjacent parameters acts by the swap on either vector (empty even value preserved).

Write \(a_n\) for the mass of arcs from a fixed left mid-edge back to the left strictly above, before any further side contact, in the infinite strip, and
\[
 {\cal B}_n=1-\tau a_n.
\]
Below gives the same mass by reversal and translation. At homogeneous honeycomb \({\cal B}_n\) is exactly the positive crossing mass from that port to the right boundary. Indeed the same convex flux identity on truncated rows gives coefficient \(\cos(3\pi/8)\) on returns to either part of the left side and 1 on crossings (turns \(\pm\pi,0\) by closing along the convex boundary). Remote top exits go to zero: they are the single-top-defect outputs after empty lower rows, one \(A_L\), then a growing number of upper source-free rows in the odd block, which is transient. Reflection bounds bottom exits. Pass to the infinite strip on left/right by positivity; \(\tau=2\cos(3\pi/8)\). We take \(a_0=0,{\cal B}_0=1\).
These strips in the physical tiling are precisely bands between triangle-side grid lines along one lattice direction (\(n\) equal spacings); the other such directions follow by symmetry.

We will need the edge reflection rule on the lower vectors. At the last column reflection \(z\mapsto1/z\) simply multiplies each component by \((U/u)^{o}\) with \(o=0,1\) its port occupation: successive contacts from the left are linked there in consecutive pairs with weights \(uU v^{\ell-1}\) (\(\ell\ge1\) their row gaps), except possibly a last contact linked to the top with factor \(u\). At the first column it multiplies an even component by \((u/U)^o\), and sends \(O\) to the following vector componentwise:
\[
 (u/U)^o O+\frac{U^2-u^2}{U(1-v)}\,D_* E.
\]
Here \(D_*\) adds a first-site defect if that port was vacant, or else deletes the first-site end, turning its partner into the defect.

For details, fix the rest of a lower-half diagram and list the first column's occupied \(R\)'s at \(t_1<\cdots<t_l\le0\), followed in this list by the final \(T\) if occupied. In the even calculation pair consecutively using the single vertical channel, with factor \(U\) for a final \(R\)-to-top link. In the odd calculation the left source attaches to an odd-index member \(j\) of this list, with consecutive pairings before and after. If it attaches to \(R_j\) and is bounded above by the next \(R\)-time \(t_+\), sum over source positions relative to the unmodified weights of all other pairs. The factor is
\[
 v+\frac{u^2+U^2}{1-v}
 +\left(w-\frac{u^2}{1-v}\right)v^{t_j-t_- -1}
 +\left(W-\frac{U^2}{1-v}\right)v^{t_+-t_j-1}
\]
where \(t_-\) is the previous time or \(-\infty\). Indeed from below/above one jumps along a free vertical segment, with a double cell at a coincident neighboring time. Both correction coefficients are \(-\gamma v\), so this respects the multiplicative reflection. For \(R_j\) last in the list, omit the final neighbor term and subtract \(U^2 v^{-t_j}/(1-v)\) (cutoff at row 0); the reflection difference is then the claimed multiple of the even \(R_j\)-to-top weight \(U v^{-t_j}\). Finally for attachment to the top itself the analogous factor is \(u/(1-v)+(w-u^2/(1-v))v^{-t_l}/u\) (omit the last term when there is no previous contact). Its reflected value minus \(u/U\) times the old gives the same stated coefficient times the vacant-top value. In the difference terms substituting the end only reroutes the defect, so loop exclusion upon gluing the rest agrees with \(D_*\). These local geometric sums prove the identities at convergent parameters and thus rationally. To justify a common open set one may start at homogeneous honeycomb. The same local formulas for a single reflected edge, summing over finite fixed rest diagrams, bound its positive sums by the original finite \(E,O\); accessibility as above gives transience also there, hence a common complex neighborhood by continuation. The one-column rules are interpreted with no intervening strip and work identically.

Use a different Pfaffian here than in the periodic problem:
\[
 H(x,y)=x^2+y^2+\sqrt2 xy-\tfrac12,\quad
 F_o=\prod_{i<j}\frac{H(x_i,x_j)}{x_i-x_j},\quad
 Q_o= F_o\operatorname{Pf}_{\rm ext}\!\left[\frac{(x_i-x_j)(x_i+x_j+d)}{H(x_i,x_j)}\right] .
\]
The augmentation is again a final column 1 in odd size. This is a symmetric polynomial of degrees at most \(n-1\), by cancellation at equal rows. At \(x_i=a\) it drops that site multiplying by \(\prod_{j\ne i}(x_j+b)\); at \(x_i\to\infty\) its leading coefficient drops the site. Indeed the skew column against other sites is then constant, supplying or eliminating the augment, using \(H(y,a)=(y-a)(y+b)\). There are recursions (remaining sites denoted by \(S\) in \(x\) coordinates):
\[
\begin{aligned}
 Q_o(j,k',S)&=(j+k'+d)\prod_{y\in S}R(y;j,k')\,Q_o(S),\\
 R(y;j,k')&=\frac{H(y,j)H(y,k')}{(y-j)(y-k')},\qquad H(j,k')=0,\\
 Q_o([q m],[m/q],S)&=\sqrt2(h+a)\prod_{y\in S}(y+h)\, Q_o(h,S),\qquad h=[m].
\end{aligned}
\]
The first is skew pole elimination (if \(j=[q^3 l],k'=[l]\), then \(R=(y-[q^6 l])(y-[q^{-3}l])\)). For the second induct in the size, interpolating in any spectator. At internal pole pairs \(H(j,k')=0\) among spectators the recursions agree by \(R([q m])R([m/q])=(h+j)(h+k')R(h)\) using the factors just given. At spectator value \(-h\) the result is zero by the first rule (simultaneous poles paired with the two sites); at \(a\) it agrees by \(([q m]+b)([m/q]+b)=(h+a)(h+b)\); at infinity the leading coefficients agree. These \(2(n-3)+2\) finite tests plus the leading term suffice for \(n\ge3\), and size two is immediate.

\subsection{Half-strip polynomials}\label{L:step7}

The next proposition supplies the denominator and degree bounds needed
when the two half strips are glued. Its proof constructs the polynomial
vectors before identifying them with the physical sums.

\Needspace{6\baselineskip}
\begin{proposition}[Polynomial half-strip vectors]\label{L:prop:strip-polynomials}
For either vector \(V=E,O\), with \(\sigma=-1,+1\) respectively, let \(o_j\)
denote the occupations of the output pattern. Then
\[
\begin{gathered}
 D_n=Q_o\prod_j z_j^{n-1},\qquad
 \psi= D_n \left(\prod_j z_j^{\sigma o_j/2}\right)V,\\
 \psi\text{ is polynomial with }\deg_{z_j}\psi\le m_j=2(n-1)+\sigma o_j.
\end{gathered}
\]
At ordered neighboring high-low parameters \(q^r l,l\) for \(r=2,3\), the vector value uses respectively \(P V\) (at fused parameter \(q l\)) and \(C V\) (dropping the pair). There are also direct wall rules:
\begin{itemize}
\item At a first parameter \(p^{-1}\), insert a vacancy on either smaller vector, and in the odd vector also insert a defect there from the smaller even vector, summing; call this \(K_L(E,O)\).
\item At a last parameter \(p\), just insert a vacancy on each (\(K_R\)).
\end{itemize}

Even vectors drop a vacant site at \(z=0,\infty\) (occupied components go to zero). Odd vectors also satisfy the following edge refinements at the first or last site \(e\). With it occupied and paired, \(V\) has vanishing leading term both at \(z_e=0,\infty\) on the scale \(z_e^{\mp1/2}\), respectively. With it the defect,
\[
 \lim_{z_e\to\infty} V/\sqrt{z_e}=c_* E_{\rm dropped},\qquad
 c_*=[z^{1/2}]g_z
\]
where the bracket on the right denotes a Laurent coefficient. All parameter rules concern generic other values.

\end{proposition}
\begin{proof}
We induct on $n$, retaining all the stated bounds, pair and wall rules,
exchange covariance, and the two row equations at smaller sizes.
There are two distinct uses of edge behavior. Vanishing at a paired
occupied edge will first identify the coordinate charts. The leading
coefficient at an edge defect will then establish exchange covariance;
only afterward will we prove the row equations at size $n$.

\paragraph{Startup and polynomial boundary data.} For startup at sizes 0 and 1 there is only empty even, and odd gives \(O_1=g_{z_1}\) at size 1. At size 2 write \(Q=Q_o(x_1,x_2)=x_1+x_2+d\), use position subscripts for weights, and label odd states by their defect. Besides the empty even we have
\[
 E_{12}=s_0 h_1 g_2/Q,\qquad
 O_1=g_1+\tau h_1(x_2-a)/Q,\qquad O_2=g_2[1-s_0(1-g_1^2)/Q].
\]
Indeed \(E_{12}=U_1u_2/(1-v_1v_2)\), and the two odd equations are
\(d_1 O_1-U_1U_2 O_2=u_1+W_1U_2 E_{12}\),
\(d_2 O_2-u_1u_2 O_1=v_1u_2+U_1 v_2 E_{12}\) (\(d_i=d_{z_i}\)); substitution using the weight identities gives the expressions. These obey the degrees, edge and wall claims. On either pair specialization the systems remain generically nonsingular (\(v_1v_2\ne1,\tau^2 v_1v_2\ne1\)); sliding the triangle/cap through the row by the diagram identities gives the pair rules.

At size \(n\ge3\) prescribe putative polynomial values \(S\) on the following loci, rescaled as in \(\psi\). For every \(h<l'\) use \(z_h=q^r z_{l'}\), \(r=2,3\) (only these linear high-before-low prescriptions). Take the triangle/cap rule made adjacent and transport back by swaps, before diagonal rescaling, moving through intervening spectators. For any site use also its left-wall value at \(z=p^{-1}\) and right-wall value at \(p\), transported likewise from the designated wall. Prescriptions are independent of the transport before partners cross each other, and covariant under such swaps: for pairs this is precisely triangle/cap sliding, braid and inversion as in the periodic construction; for a wall one inserts at the edge into the smaller ordered list, and swaps of other positions commute past the insertion (and past its transports by braid).

Each pair or wall value has polynomial dependence of degree \(\le m_j\) in spectators \(z_j\). Indeed the smaller denominator cancels by the \(Q_o\) rules. For an intervening site between partners one can choose to cross either partner in transporting. Writing the low parameter as \(l\), the common possible swap poles are \(z_j=q^5 l\) in fusion (killed by \(x_j+[q l]\)), and \(q^5 l,q^6 l\) in the other case (killed by \(R\)). For wall transports the poles are \(q^{7/2},q^{9/2}\), killed by \(x_j+b\). Poles at these swaps are at most simple at a generic point. Integral powers follow by port parity (coherent roots on a pair, i.e. \((q^r l)^{1/2}=q^{r/2}l^{1/2}\)). 

At a spectator extreme, put \(M=\max(|z_j|,|z_j|^{-1})\). Before the \(\psi\) rescaling the values are bounded by \(O(M^{\sigma o_j/2})\): swaps crossing this site are bounded and at leading order preserve its spectral-label occupation (opposite weight \(v\to1\), both-square doubles bounded, single corners \(O(M^{-1/2})\)). Wall insertion from even into odd is harmless. This gives the stated bounds.

At a pair's simultaneous extremes, writing \(z_i\asymp z_j\) for its two parameters in order of magnitude, the common leading scalar behavior is
\[
 Q_o\ \sim\ (x_i+x_j)(x_i x_j)^{n-2} Q_{o,{\rm rest}}
\]
by the two recursions, so \(D_n\) has order 1 at the pair origin and growth order \(4n-5\) at pair infinity. The prescribed vector before rescaling is bounded if \(o_i+o_j=0,2\); if \(o_i+o_j=1\) it is \(O(M^{\sigma/2})\). Indeed fusion with one leg occupied uses the smaller occupied site with that bound, and changes of occupation count parity in transporting cost a small single corner. Thus the rescaled \(S\) along the relation has order \(\ge1\) at zero and degree \(\le m_i+m_j-1\) at infinity, except for both occupied and even, where the bounds are 0 and \(m_i+m_j\). 

In the exception the values at zero, and leading coefficients in coordinates \(S/(z_i^{m_i}z_j^{m_j})\) at infinity, agree between the two gap prescriptions. For the vector only the local cup (\(X\) or 1) contributes there before transporting, because the fused smaller vector at a vacant port drops it and the occupied case vanishes. Leading swaps preserve occupations, with limiting coefficients independent of gap. Indeed using the displayed leading scalar in the diagonal rescaling at \(z_i/z_j=\lambda=q^{\pm r}\) gives a factor proportional to \(\lambda^{1/2}+\lambda^{-1/2}\) in either limit in question (the proportionality factor independent of gap). Including the local cup gives the same multiplier since \(X(q+q^{-1})=q^{3/2}+q^{-3/2}\). For both occupied and odd we also note the sharper bounds 2 and \(m_i+m_j-2\).

\paragraph{Compatibility of pair and wall data.}
The prescribed values are polynomial with the required spectator bounds.
To interpolate them simultaneously, we now check their agreement at generic
intersections of different constrained unordered pairs or sites:
\begin{itemize}
\item For pair-pair checks the same linear-order arguments as for periodic interpolation work. Disjoint reductions commute after regular swaps preserving their orientations. At a shared high or shared low site, identical gaps give identity exchange on the equal other parameters; unequal ones can be sorted into consecutive exponents \(3,1,0\) or \(3,2,0\) (relative to a shift). Two triangle insertions then equal the cap insertion at gap three with a swap moving the middle back between: a cell \(B(q)\) with leg continued by the cap is just the two triangles. Consecutive gaps both two give triangle then cap either way (empty or any cup of weight \(X\)). Other consecutive gaps give the wheel exponents (three circular gaps \(2,3,3\) in some order); both polynomial values vanish by \(x_j+[q l]\) or \(R\) in the reduction ratio. In this wheel check one brings each constrained pair together without crossing the third site (linear orders), so no pole cancels the zero.
\item A wall and disjoint reduction commute by rearranging to direct insertions with regular swaps. Equal-side wall constraints on distinct sites likewise agree (equal parameters with identity exchange); left before right commute directly. The reverse order of opposite walls is also a gap-three pair considered below.
\item For a pair and a wall on its high site at \(p^{-1}\), or low site at \(p\), both values vanish: the other site has \(x=-b\), and the fused \(x\) is \(-a\) or the pole-pair sum factor vanishes. The wall need not cross its partner here.
\item For a fusion pair with the \textbf{low} site \(p^{-1}\), make the pair first. The wall rule adds \(K_L\) before the high site \(q^{1/2}\), and then swaps them by \(B(q^2)=P P^t\). Now \(P^t K_L\) gives the reflected smaller value (at \(q^{-1/2}\)): it multiplies occupied components by \(X=u/U\) there (ratio taken at the old high site), and in the odd vector also adds \(X D_* E\), exactly the edge reflection proved above. Symmetrically with the high site \(p\) make the pair last; \(P^t K_R\) reflects the low site \(q^{-1/2}\) by its factor \(X=U/u\) on occupied ports.
\item In the corresponding gap-three case the ordered pair is \(p,p^{-1}\). Make it first. The low-site wall adds \(K_L\) before the high site and swaps back by \(CC^t\). To compute \(C^t K_L\) use induction on the high site of the smaller vector: it comes as a vacancy from the right wall moved to the left through the remaining sites by \(B(p/z_j)\). These swaps are exactly an ordinary row, with right end vacant and left leg brought out as that high site: the swap ports are now \((B,R,T,L)\), a quarter turn, so its weights become those of the row at \(p z_j=q^3/(p/z_j)\). Contracting both sites vacant gives \(T E,T O\); contracting both occupied, low the defect from \(K_L\) on even, attaches it to the left leg and gives \(A_L E\). Thus one gets precisely the smaller row invariance. For the high site's right-wall value, instead make the pair last and use the low site's smaller insertion from the left through swaps \(B(z_j/p^{-1})\), now just \((L,B,R,T)\). Contracting by \(C^t K_R\) requires vacant outgoing auxiliary; the incoming one came from \(K_L\), giving again \(T E,T O+A_L E\). Both wall checks therefore agree with the cap rule.
\end{itemize}
Here reordering with arbitrary other variables is before diagonal rescaling using covariance of prescriptions. One can choose coherent branches in the nonvanishing checks. At vanishing denominators the zero arguments use the \emph{ratio} after the smaller denominator is already canceled.

\paragraph{Chart construction and the paired-edge bound.}
For a pattern and a distinguished site \(i\), form a chart by interpolation in \(z_i\) using both gaps with all other positions in the designated linear order. Add no wall if even occupied at \(i\); one wall (either) if vacant; both if odd occupied. These are \(m_i+1\) nodes. 

Each chart is polynomial in all sites with the claimed degrees. Indeed nonzero node coincidences yield only simple poles canceled by compatibility. When \(z_j\to0,\infty\), the two moving nodes from \(j\) have Lagrange factors of orders \(O(|z_j|^{-1})\) at zero and \(O(|z_j|^{1-m_i})\) at infinity, covered by the pair estimates; the fixed nodes use just the spectator bounds. In the even both-occupied exception the potential pole cancels between the two nodes thanks to their equal origin limits; the potential degree \(m_j+1\) cancels by the matched leading coefficients. Explicitly if their locations are \(\lambda z_j,\lambda' z_j\), the offending two infinity multipliers after inserting \((\lambda z_j)^{m_i}\) or \((\lambda'z_j)^{m_i}\) are proportional to \(-\lambda'\lambda/(\lambda-\lambda')\) and its exchange.

Every chart also obeys pair prescriptions not involving \(i\): check at its interpolation nodes by compatibility. None are forced to coincide on an ordered pair relation (in particular if \(i\) is between the partners the two node exponents cannot collide since a prescribed gap is not a sum of two allowed gaps modulo 8). Likewise it obeys a wall at a different position when no nodes are forced to coincide there, in particular the direct near wall (\(p^{-1}\) first or \(p\) last).

For odd and an edge \(e\ne i\) occupied and paired, the chart has vanishing top and constant coefficients in \(z_e\). In fixed-node terms (nodes not from \(e\)), only a smaller defect at \(e\) can contribute vector growth of order \(M^{1/2}\) there by induction. Pair transports can avoid \(e\); if a wall transport crosses it, the added site from this leading term is vacant when crossing (first step from the edge), so the defect remains at \(e\) directly, not paired. 

For moving nodes from \(e\), if \(i\) is occupied the sharper odd bound for both occupied already suffices. If \(i\) is vacant, make the pair adjacent at the edge \(e\) and then move \(i\) back. Leading growth \(M^{1/2}\) before rescaling requires the fused site occupied as defect by induction. Its link must go to \(e\) (otherwise \(e\) stays vacant), which then stays defect, so this term too cannot contribute.

Choose as guide a first or last position \(e\) not the defect. All charts \(i\ne e\) for the pattern coincide: in variable \(z_e\) they obey the pair nodes, plus the direct near wall if vacant. This suffices by degree (in the odd occupied case divide by \(z_e\) using the two vanishing coefficients). If both edges are guides their results agree by an interior chart. 

Denote the common result from interior charts, undone to vector scaling, by \(\widehat E,\widehat O\). Thus they obey all pair prescriptions, both direct near walls, and both transported walls at any interior site (for a vacant or odd occupied site use its own chart; in even patterns where needed test the wall by another chart avoiding a collision, e.g. one on the later side for \(p^{-1}\), earlier for \(p\)). The even empty value is 1 since empty evaluation just sees empty in all prescriptions (loops killed). The paired-edge bound just proved holds by taking interior charts.

We have now constructed $\widehat E,\widehat O$ with the component degree
bounds, all pair prescriptions, the direct near-wall rules and the
interior transported-wall rules. We also know the improved bound at a
paired occupied edge. These conclusions concern the interpolated vectors;
row invariance and identification with the actual sums are still to come.

\paragraph{The leading coefficient at an edge defect.}
We prove the remaining leading edge claim at infinity for \(\widehat O\)
with the defect at \(e\). In polynomial coordinates its desired coefficient of \(z_e^{m_e}\) is
\[
 L_0= c_*\,2^{-(n-1)}\left(\prod_{j\ne e} z_j^{1+o_j}\right)
       \psi_{\rm even,dropped}.
\]
It has degree at most \(2n-3\) at an interior position \(i\). The actual leading coefficient \(L\) agrees at the fixed nodes of chart \(i\) (all except the two from \(e\)) by induction: in pair prescriptions use transports not touching the edge, and in walls crossing \(e\) the added port from the old leading defect is vacant and swaps past by straight spectral-label continuation, exactly giving the even wall rule on dropping \(e\). If \(i\) is occupied, moving terms cannot contribute to \(L\), so its degree is at most \(m_i-2=2n-3\) from the remaining Lagrange basis; there are \(2n-2\) fixed tests. If \(i\) is vacant, the moving terms by the pair bounds yield degree at most \(m_i-1=2n-3\) in \(L\) (their leading basis factors just use the product over fixed nodes); here there are \(2n-3\) fixed tests. It suffices to match the highest coefficient \(A_0=[z_i^{m_i-1}] L_0\).

Indeed in this vacant-\(i\) case along \(z_e=y\to\infty,z_i=\lambda y\), the polynomial divided by \(y^{m_e}(\lambda y)^{m_i-1}\) tends to \(A+B'\lambda\), with \(A\) that coefficient in \(L\). Use \(\lambda_r=q^{\pm r}\) with fixed sign according to the partner order. At gap three the limit vanishes by the cap bound (no growing defect). At gap two, fusing at the edge, before rescaling the vector divided by \(\sqrt y\) tends to \(c_* X\lambda_2^{1/4} E_{{\rm drop}\ i,e}\) by induction (the fused defect leads to \(e\); \(i\) vacant then transports back by straight continuation). Using the simultaneous leading \(Q_o\) above, this gives
\[
 A+B'\lambda_2=(1+\lambda_2) X\lambda_2^{1/4} A_0 ,
\]
where in \(A_0\) the smaller even vacancy at \(i\) itself drops at infinity. All fourth-root expressions here use the signed exponent before taking the root, according to the same partner order. Since \(X\lambda_2^{1/4}(1+\lambda_2)=1-\lambda_2/\lambda_3\) and \(A+B'\lambda_3=0\), we have \(A=A_0\), proving equality with \(L_0\).

\paragraph{Exchange covariance and all wall rules.}
Both edge estimates are now available. They provide the final interpolation
test in the only case where the pair and wall nodes alone do not determine
an exchange relation. The constructed vectors therefore have swap covariance. For an adjacent swap, test the relation in a third variable \(z_j\) at the pair nodes involving \(j\), valid by covariance of prescriptions, and its available walls. Both sides have the polynomial degree \(m_j\) after the indicated spectator rescaling since the swap does not touch that occupation. Only an odd occupied edge with just its near wall could need one extra test; at infinity the leading vector there is zero when paired and \(c_*\) times the smaller even when defect, already covariant under this swap (which cannot turn the directly carried defect into a pair). This suffices. By transport we now have \emph{all} wall prescriptions.

\paragraph{Row invariance and physical identification.}
It remains to show that the constructed vectors solve the two row equations.
Apply row addition to these vectors and subtract them (including \(A_L\widehat E\) in the odd row). Denote the residual by \({\cal R}\). After multiplication by \(D_n\prod_j G_j\prod z_j^{\sigma o_j/2}\), \(G_j=z_j^2+2b z_j+1\), it is polynomial of degrees \(\le 2n+\sigma o_j\): the row denominators are cleared and a change of bottom/top occupation costs a single small corner at extremes. It vanishes on all pair and both wall loci, by sliding triangle/cap and swaps and using smaller row invariance. At a direct first wall the cell has only \(LT,BR\) routes (weight 1); an old inserted defect goes to the smaller row's left leg and a new source from \(A_L\) goes to the top, giving invariance of \(K_L\). At the last wall \(LB,TR\) routes just carry vacancy as required.

Thus with
\[
 H_+=\prod_{i<j}(z_i-q^2 z_j)(z_i-q^3 z_j),\qquad
 H_\partial=\prod_j(z_j-p)(z_j-p^{-1}),\qquad
 \widetilde{\cal R}=\frac{D_n\prod_j G_j}{H_+ H_\partial}\,{\cal R},
\]
the even residual vanishes (empty is immediate; otherwise negative polynomial quotient degree). Each odd \(\widetilde{\cal R}\) is independent of vacant variables and supported on powers \(\pm1/2\) at each occupied one. Moreover it vanishes for a paired occupied edge: the paired improved bound at the two extremes survives row addition there. Indeed in the source-free odd row the only leading way to carry an occupied edge to itself is vertical propagation (side end forbidden); an old defect then stays defect; \(A_L\widehat E\) is bounded. The scalar quotient here is itself bounded at a generic single-site extreme.

Swapping neighbors sends the old \(\widetilde{\cal R}\) by the swap to the new multiplied by \(H_+(\text{new})/H_+(\text{old})\). For exactly one occupied site among them, its equation couples just the two slid patterns, with single-route coefficients \(u,v\); vanishing of one forces vanishing of the other. In every pattern with at least three occupied sites a leftmost or rightmost occupied end is paired and can slide via vacancies to the corresponding edge. Thus all such residuals vanish. For a lone defect taken at the edge, there is a swap on two vacancies (\(n\ge3\)); it acts there just as identity now, while the indicated multiplier is nonconstant and the component independent of both parameters. This gives zero too, and one slides to all sites.

Thus
\[
 T\widehat E=\widehat E,\qquad
 T\widehat O+A_L\widehat E=\widehat O.
\]
Together with $\widehat E_\emptyset=1$, uniqueness near homogeneous
honeycomb parameters identifies $\widehat E,\widehat O$ with the actual
rational vectors $E,O$. The even vacancy extreme assertion used inductively now follows at this size by taking row-equation limits there: occupied components vanish by the degree bound, and with bottom/top vacant the row cell just carries the auxiliary line unchanged in the limit. Smaller row uniqueness and the empty value give removal. This finishes the induction.
\end{proof}


\subsection{A crossing determinant}\label{L:step8}

To form \(a_n\), glue \(A_L E\) in the fixed source row to the upper reflected \(O\), matching occupations and requiring precisely the single path, no loops. This proves rational continuation. Sliding spatial swaps gives symmetry (opposite ratios in the reflected vector; diagram transpose of a swap has the same weights, then use inversion). The same gluing reduces at a gap-three pair to the remaining quantity without multiplier: carry the cap from below through the source row, then contract by \(C^t\) above. This gives the smaller reflected vacuum there since its exchange on the reversed rapidity relation uses \(CC^t\), the exchanged value itself uses \(C\), and \(C\) is injective. At \(z_n=p\) lower and row carry a last vacancy; the vacant-last component above also drops by right reflection followed by right-wall removal. Finally \(a_n\) is invariant under \(z_n\mapsto1/z_n\) by the full-strip arc sum: all occupations of that last column are consecutive paired contacts from the left, with product weights involving \(uU,v\) only. As for edge reflection before, both sides converge on a common open set (upper halves handled by vertical reflection). Symmetry gives separate inversions everywhere.

Consequently
\[
 P_n={\cal B}_n Q_o^2\prod_j(x_j+b)
\]
is a symmetric polynomial in \(x\) of degrees \(\le2n-1\). The half-strip polynomial bounds and the inserted row denominator make it Laurent in \(z\) (half-powers match at occupied ports); inversion then gives polynomiality in \(x\). Before rescaling the gluing stays bounded at single extremes: if the upper odd occupation grows it faces either a lower even occupied end decaying or a decaying single corner. Both the wall specialization \(x=a\) and pair specialization \(H(x_i,x_j)=0\) reduce \({\cal B}\) without multipliers.

Define
\[
\begin{aligned}
 \Delta(x,y)&=2x+\sqrt2 y+s_0,\\
 N(x,y)&=(a+x+\sqrt2 y)\Delta(x,y)-2H(x,y)\\
       &=\sqrt2 xy+(2a+s_0)x+\sqrt2(a+s_0)y+\tfrac32 ,\\
 M_o(x,y)&=\frac{x-a}{x+b}\frac{N(x,y)}{H(x,y)}.
\end{aligned}
\]
Then
\[
 {\cal B}_n=\frac{F_o^2}{Q_o^2}\det M_o(x_i,x_j).
\]
Here and initially in what follows statements using kernels are rational before specialization. Indeed the proposed numerator \(P_n\) is polynomial with the indicated degree: \(H\) poles cancel, including on the diagonal where \(M_o(x,x)=1-\tau^2 v\) by substitution, coincident positions cancel twice by rows and columns, and at a single infinity the site's two off-diagonal directions decouple and its diagonal tends to \(\gamma\). Thus the leading \(x^{2n-1}\) coefficient is \(\gamma\) times the smaller proposed numerator. At \(a\) the off-diagonal row vanishes and the diagonal is 1, giving wall reduction. At \(H(x,y)=0\), the double-pole determinant term gives pair reduction by
\[
 -(x-a)(y-a) N(x,y)N(y,x)=(x+b)(y+b)(x-y)^2(x+y+d)^2 .
\]
Indeed modulo \(H\), \((x-a)\Delta=(x-y)(x+y+d)\) and \((a+x+\sqrt2 y)(a+y+\sqrt2 x)=(x+b)(y+b)\). Thus by induction from size zero the only ambiguity at each size in \(P_n\) is a constant multiple of \(\prod_{i<j}H(x_i,x_j)\prod_i(x_i-a)\).

We fix it at infinity in the last column with all other columns regular honeycomb. Vary \(v\) there from its honeycomb value towards 1. Every visit to that column gives consecutive arches from its left edge of weights \(K_\ell=uU v^{\ell-1}=\tau(1-v)v^\ell\) (\(\ell\ge1\)), total \(\tau v\). The excess left-return mass above \(a_{n-1}\) consists of trains of \(k'\ge1\) disjoint ordered upward jumps with these weights. Inside the old strip they use one initial bridge from the source to the interface, one ending bridge to the left boundary, and \(k'-1\) intermediate upward interface-to-interface excursions, all mutually avoiding. Indeed for \(2k'\) successive interface ends, the outer arches pair \((1,2),(3,4),\ldots\). The two inner bridges have terminals in ascending order by planarity; no inner pair spans either terminal. Any nonempty prefix/suffix outside them would thus close into loops, so they end at \(1,2k'\). Then 2 must pair inside to 3 (a later partner again traps loops), and so on. Conversely these pieces in order and without conflicts form precisely such arcs if the last bridge ends above the source. For width zero of the old strip use the empty bridges only. By half-turn symmetry, translation and reversal the final bridge and intervening excursion masses per start are \({\cal B}_{n-1}\) and \(a_{n-1}\).

Hence
\[
 0\le \Delta a_n\le
 \frac{{\cal B}_{n-1}^2\tau v}{1-\tau v a_{n-1}},
 \qquad \lim_{v\uparrow1}\Delta a_n=\tau {\cal B}_{n-1}.
\]
Here \(\tau a_{n-1}<1\). For equality in the limit, at fixed train length normalize each jump by its total: the gaps diverge in probability. For any fixed shapes of the finite interior pieces this eventually removes all avoidance and end constraints with probability tending to 1. Dominated convergence and the geometric train tail suffice. This calculation defines an analytic summed arc mass along the interval \(v<1\) in question (locally, \(\sum |K_\ell|a_{n-1}<1\)) and agrees with the actual gluing near honeycomb. Thus the polynomial identity defining \(P_n\) using that mass continues along it even with fixed coincident other parameters or a zero of \(Q_o\). Since \({\cal B}_n\to\gamma{\cal B}_{n-1}\), the leading coefficient there agrees with the determinant's. The possible ambiguity factor has nonzero leading coefficient there (\(x_j=b\) for others), proving the formula.

It suffices to estimate even sizes at \(x_j=1\) (square rhombi), by
\[
 {\cal B}_{2n}(b,\ldots,b)\ \le\ {\cal B}_n(1,\ldots,1)\ \le\ {\cal B}_n(b,\ldots,b).
\]
Here are convergence and comparison details. In mixtures of the two angles \(\pi/3,\pi/2\) all weights are nonnegative with single routes strictly positive and doubles bounded by the product of the two single weights (\(w\le u^2,\ W\le U^2\); use the explicit weights/identities). One last honeycomb column is arch-by-arch dominated by one square since \(0<v_H<v_S<1/2\) in \(K_\ell\). Conversely two honeycomb columns dominate one square there: for each arch interval restrict to paths going up each row carrying one occupied vertical channel, giving mass at least
\[
 uU(1,v)
 \begin{pmatrix}v&u^2\\ U^2&v\end{pmatrix}^{\ell-1}
 \binom{1}{v}
\]
with honeycomb \(u=X,\ U=v=X^2\); the two off-diagonals switch channels via two corners, and the end factors allow starting/ending via a horizontal propagation. The matrix sends the indicated column vector to at least \(v(1+u^2)\) times itself, \(>v_S\) in multiplier, and \(uU(1+v^2)>u_S U_S\). For example \(X>.541,\ .2928<v<.294,\ .320<v_S<.322,\ 2a<.766\) by the radicals above, \(u_SU_S=2a(1-v_S)v_S\). Distinct intervals use disjoint rows and have the required connectivity.

Induct on the number of squares simultaneously over widths to get transience and convergence. Spectra on the nonempty even and odd blocks are unchanged by permuting (nonsingular swaps at these ratios), also for vertically reflected angles. With a square placed last the two-column domination bounds \(a_n\) by a finite one by induction. This also bounds all lower and upper arc components: drop mutual conflicts of separate arcs using the product bound (self double visits keep their prescribed weights). For a single lower arc between fixed top ports append disjoint fixed tails in fresh rows to two left ports (route the nearer one first horizontally to the wall continuing the farther vertically, then the other); for a variable left-to-top arc append one such fixed tail. The resulting masses are bounded by full left-return sums by translation and reversal. Reflection interchanges the two directions of return and gives the upper version. Accessibility gives transience as before, hence also in every permutation, where gluing again gives the actual finite mass. We may thus use swap symmetry and repeatedly move a column to the last wall to apply either comparison, proving the displayed bounds. The honeycomb value is also nonincreasing in width by return inclusion.

\subsection{Square-point asymptotic}\label{L:step9}

We now estimate the crossing determinant at $x_1=\cdots=x_n=1$,
for even $n$. The positive column comparison proved above will then
transfer the estimate to every honeycomb width. There are two reductions.
First we express the large determinant quotient as a quotient of residual
Hermite evaluation determinants of sizes three and one. We then evaluate
the cancellation in the three-column determinant; its first nonzero
normalized term has order $n^{-3}$.

\paragraph{From the crossing kernel to interpolation remainders.}
For even $n$, adding the all-ones matrix to the skew kernel defining
$Q_o$ does not change its determinant: the adjugate of an even-size skew
matrix is skew, and its quadratic form on the all-ones vector is zero.
The new numerator is
\[
 H(x,y)+(x-y)(x+y+d)=(x-a)\Delta(x,y).
\]
Consequently the crossing determinant gives the rational identity
\[
 {\cal B}_n
 =\prod_i\frac1{x_i+b}\,
   \frac{\det[N(x_i,x_j)/H(x_i,x_j)]}
        {\det[\Delta(x_i,x_j)/H(x_i,x_j)]}
 =\frac{\det K_{S_3}(x_i,x_j)}{\det K_{S_1}(x_i,x_j)}.
\]
Here $K(x,y)=1/H(x,y)$, and $K_S(x,y)$ is $K(x,y)$ minus its
polynomial interpolant in $y$ of degree less than $|S|$ at the Hermite
nodes
\[
 S_1=(b),\qquad S_3=(b,b,-a).
\]
The repeated node in $S_3$ prescribes both value and first derivative.
To verify the second equality, use the polynomial congruences
\[
\begin{split}
 (y-b)\Delta(x,y)&\equiv r(x)=-\sqrt2(x+a)(x+b),\\
 (y-b)^2(y+a)N(x,y)&\equiv(x-b)(x+b)(x+a)r(x)
 \pmod{H(x,y)}.
\end{split}
\]
The second follows from $N\equiv(a+x+\sqrt2y)\Delta$ and
$(y-b)(y+a)\equiv(x+a)(a-x-\sqrt2y)$ modulo $H$.
Multiplying either numerator by its node polynomial and dividing by the
indicated row factor produces $K$ plus a polynomial of the required
interpolation degree. Its Hermite data vanish, so it is exactly $K_S$.
The row and column factors cancel to give the displayed quotient.
All these identities are rational identities before confluence; we will
prove that the denominator is nonzero at the square point.

\paragraph{Confluence as a positive Gram matrix.}
Use the positive measure and its orthonormal polynomials from
Section~\ref{L:part4}:
\[
 c=\frac34,\qquad
 w_b(\nu)=\frac{k}{\sqrt2}\frac{\sinh(\pi\nu/4)}{\sinh(\pi\nu)},
 \qquad
 f_j(\theta)=\int p_j(\nu)e^{\theta\nu}w_b(\nu)\,d\nu.
\]
The leading coefficient of each $p_j$ is positive. The calibrated Jacobi coefficients and Lemma~\ref{L:lem:polynomial-laplace} give
\[
 b_j^{J}=\frac{j}{2c}
       \left(1+\frac{d_J}{j^2}+o(j^{-2})\right),
 \qquad
 f_j(\theta)=C(\theta)\rho(\theta)^j
       \left(1+\frac{d(\theta)}j+o(j^{-1})\right),
 \quad
 \rho(\theta)=\tan\frac{\theta}{2c},
\]
for fixed $0<\theta<c\pi$, with $C(\theta)>0$; the Jacobi diagonal is
zero. The convergent product $S_j$ in the earlier Laplace estimate has
been absorbed into this first-order expansion.

The useful Gram representation is
\[
 K(x,y)=\frac2k\int V_\nu(x)V_\nu(y)w_b(\nu)\,d\nu,
 \qquad
 V_\nu(\cosh s)=\frac{\sin(\nu s)}{\sinh s}.
\]
Indeed,
\[
 H(\cosh s,\cosh t)
 =\left(\cosh(s+t)+\frac1{\sqrt2}\right)
  \left(\cosh(s-t)+\frac1{\sqrt2}\right),
\]
and the formula follows by subtracting the Fourier transforms at $s-t$
and $s+t$. Convergence and analytic continuation cover the arguments
used below. Put
\[
 q_r(\nu)=\frac1{r!}\partial_x^r V_\nu(x)\big|_{x=1},
 \qquad 0\le r<n.
\]
These polynomials have odd degrees $1,3,\ldots,2n-1$ with nonzero
leading coefficients. Thus their span is the full odd polynomial space
of degree at most $2n-1$, and
\[
 G_{rs}=\frac2k\int q_r(\nu)q_s(\nu)w_b(\nu)\,d\nu
\]
is positive definite. Dividing the determinant of $K(x_i,x_j)$ by its
two Vandermonde factors and sending all $x_i$ to $1$ gives $\det G$.
The same confluence applied to $K_S$ is a finite-rank perturbation of $G$.

Here is that perturbation explicitly. Write $d_S=|S|$, and let
$\lambda_\ell$, $1\le \ell\le d_S$, be the Hermite functionals in the
$y$ variable: evaluation at $b$ for $S_1$, and evaluation at $b$,
first derivative at $b$, and evaluation at $-a$ for $S_3$.
Use the polynomial interpolation basis
\[
 e_m(y)=V_{im}(y),\qquad 1\le m\le d_S.
\]
The polynomial $e_m$ has degree $m-1$ and nonzero leading coefficient,
so the matrix $E_{\ell m}=\lambda_\ell(e_m)$ is invertible. Define also
\[
 B_{rm}=q_r(im),\qquad
 C_{r\ell}=\lambda_\ell^y\left[
           \frac1{r!}\partial_x^r K(x,y)\big|_{x=1}\right].
\]
The jet matrix of the subtracted interpolant is $CE^{-\mathsf T}B^{\mathsf T}$.
The finite-rank determinant identity therefore gives
\[
 \frac{\det(G-CE^{-\mathsf T}B^{\mathsf T})}{\det G}
 =\frac{\det(E-C^{\mathsf T}G^{-1}B)}{\det E}.
\]
This displays the promised reduction from size $n$ to size $d_S$.
For a function $f$ with finite polynomial moments, define its moment
projection by
\[
 \Pi_n f=\sum_{j=1,3,\ldots,2n-1}p_j\int p_j(\nu)f(\nu)w_b(\nu)\,d\nu.
\]
This agrees with orthogonal projection whenever $f$ belongs to $L^2(w_b)$,
but needs only the displayed moments. The small residual matrix has entries
\[
 \big(E-C^{\mathsf T}G^{-1}B\big)_{\ell m}
 =\lambda_\ell^y\left[V_{im}(y)
             -(\Pi_n V_{\bullet}(y))(im)\right].
\]
The moment integrals are finite at all the Hermite data used here.
Only polynomial evaluation is used at $im$; no square-integrable test
vector at that argument is required.

Set
\[
 \alpha=\frac\pi8,\qquad \beta=\frac{5\pi}8.
\]
For $y=\cos\theta$, the projection in the last display becomes
\[
 V_{im}(\cos\theta)
       -(\Pi_n V_{\bullet}(\cos\theta))(im)
 =\frac{\mathscr R_n(\theta,m)}{\sin\theta},
 \qquad
 \mathscr R_n(\theta,m)=i\sin(m\theta)
       -\sum_{j=1,3,\ldots,2n-1}p_j(im)f_j(\theta).
\]
Multiplying the value rows by $\sin\theta$ and combining the derivative
row with the value row at $\alpha$ converts the three Hermite rows into
value at $\alpha$, $\theta$ derivative at $\alpha$, and value at $\beta$.
These operations and $\det E$ involve only fixed nonzero constants.
Thus it remains to estimate
\[
 \frac{\mathscr D_n}{\mathscr R_n(\alpha,1)},\qquad
 \mathscr D_n=
 \det\begin{pmatrix}
 \mathscr R_n(\alpha,1)&\mathscr R_n(\alpha,2)&\mathscr R_n(\alpha,3)\\
 \partial_\theta \mathscr R_n(\alpha,1)&\partial_\theta \mathscr R_n(\alpha,2)&
                  \partial_\theta \mathscr R_n(\alpha,3)\\
 \mathscr R_n(\beta,1)&\mathscr R_n(\beta,2)&\mathscr R_n(\beta,3)
 \end{pmatrix}.
\]

\paragraph{Imaginary evaluations and the cancellation scale.}
Put
\[
 L=2n+1,\qquad A_j(m)=i^{-j}p_j(im)>0,
 \qquad h_m=cm-\frac12 \quad(m=1,2,3).
\]
We will normalize each residual by $i^L A_L(m)f_L(\theta)$.
All three normalized columns have the same leading term. Their first
independent differences have orders $L^{-1}$ and $L^{-2}$, producing
the determinant factor $L^{-3}$. For that reason we need two finite
differences of $A_j(m)$, rather than just its leading power.

The Jacobi recurrence gives
\[
 b_{j+1}^{J}A_{j+1}=mA_j+b_j^{J}A_{j-1}.
\]
Over two steps the positive pair $(A_j,A_{j+1})$ is multiplied by
\[
 I+\frac1j\begin{pmatrix}-1&2cm\\2cm&-1\end{pmatrix}+O(j^{-2}).
\]
Its sum therefore grows as $j^{h_m}$ with a positive limiting
coefficient: after removing this power the relative errors are summable.
The difference divided by the sum contracts by
$1-4cm/j+O(j^{-2})$, with $O(j^{-2})$ forcing, and tends to zero.
This also identifies the same leading coefficient in the two parities.
In particular,
\[
 A_j(m)\sim c_m'j^{h_m},\qquad c_m'>0.
\]
For $\Delta_2 A_j=A_{j+2}-A_j$, write the exact recurrence as
\[
 \Delta_2 A_j=e_jA_j+t_jA_{j+1},\qquad
 e_j=\frac{b_{j+1}^{J}}{b_{j+2}^{J}}-1
       =-\frac1{j+2}+o(j^{-2}),\qquad
 t_j=\frac m{b_{j+2}^{J}}=\frac{2cm}{j+2}+O(j^{-3}).
\]
It follows that $j\Delta_2A_j/A_j\to2h_m$.
Differencing this exact expression once more, using
$j^2\Delta_2 e_j\to2$ and $j^2\Delta_2t_j\to-4cm$, gives
\[
 j^2\frac{\Delta_2^2 A_j}{A_j}\longrightarrow4h_m(h_m-1).
\]
These limits also give uniform first- and second-difference bounds at
all sufficiently large indices.

Define
\[
 U_m(\theta)=\frac{\mathscr R_n(\theta,m)}{i^L A_L(m)f_L(\theta)}.
\]
For $\theta=\alpha$, where $\rho<1$, the full odd polynomial series
reproduces $i\sin(m\theta)$, so this residual is the forward tail
$j=L,L+2,\ldots$. To justify reproduction, compare Taylor coefficients:
each odd monomial has its exact finite polynomial expansion. The
estimates above give local convergence for $\rho<1$, and the
nonnegative Taylor coefficients of $f_j$ supplied by the positive
Jacobi recurrence justify absolute Taylor interchange. For
$\theta=\beta$, where $\rho>1$, use the negative finite prefix
backwards; the free term is exponentially negligible after normalization.
Thus $U_m$ is the signed sum of
\[
 \frac{A_{L+d}(m)}{A_L(m)}\frac{f_{L+d}(\theta)}{f_L(\theta)}
\]
with coefficient $(-1)^{d/2}$ for even $d\ge0$ at $\alpha$, and
$-(-1)^{d/2}$ for even $d<0$ at $\beta$, apart from that negligible term.
The forward and negative-backward geometric sums both give
\[
 S(\rho)=\frac1{1+\rho^2}.
\]

The normalization of the derivative row is legitimate: subtract
$f_L'(\alpha)/f_L(\alpha)$ times the first row before dividing it by
$f_L(\alpha)$. The normalized determinant therefore has rows
$U_m(\alpha),U_m'(\alpha),U_m(\beta)$. Set
\[
 \ell_m=\frac{\Delta_2 A_L(m)}{2A_L(m)},\qquad
 \eta_L=\frac{\ell_3-\ell_1}{\ell_2-\ell_1}.
\]
Since $L\ell_m\to h_m$ and the $h_m$ are distinct, this is defined
for all sufficiently large $L$. Keep the first column, replace the
second by $L(U_2-U_1)$, and replace the third by
\[
 L^2\big[U_3-U_1-\eta_L(U_2-U_1)\big].
\]
The first two terms in the discrete Taylor expansion
\[
 \frac{A_{L+d}(m)}{A_L(m)}
 =1+d\ell_m+
  L^{-2}\left(\frac{h_m(h_m-1)}2d(d-2)+o(1)\right)
\]
are annihilated exactly in the third column. The remaining coefficient
has a nonzero limit because $h(h-1)$ is quadratic and
$h_1,h_2,h_3$ are distinct. The second column has a nonzero linear
coefficient. Hence, up to fixed nonzero column factors, the limits of
the three columns are the same Hermite rows of
\[
 S,\qquad DS,\qquad(D^2-2D)S,\qquad D=\rho\partial_\rho.
\]
The operator $D$ inserts the exponent $d$ in either geometric sum.

We include the bounds needed to pass to these limits. After subtraction
of the constant and linear terms, the discrete Taylor remainder
multiplied by $L^2$ is bounded by $C(1+|d|)^2$ when
$L/2\le L+d\le2L$, by the second-difference bound. For fixed $d$,
$f_{L+d}/f_L\to\rho^d$. Throughout that central range the same
ratios are bounded by geometric factors times a polynomial, and the
remaining indices contribute exponentially small tails, even after
multiplication by the displayed powers of $L$.
The derivative row has the same control. Indeed the Jacobi recurrence
and the first-order Laplace expansion give
\[
 \frac{f_j'(\theta)}{f_j(\theta)}
 =\frac{j(\rho+\rho^{-1})}{2c}+\frac\rho{2c}+o(1).
\]
Subtracting the corresponding expression at $L$ proves convergence of
the derivative of $f_{L+d}/f_L$ to the derivative of $\rho^d$ for fixed
$d$, and bounds it by the same geometric factor times $C(1+|d|)$ in
the central range. Polynomial factors do not affect the exponential
tail bounds. This proves convergence of all three Hermite rows.

\paragraph{Nonvanishing and the physical exponent.}
The three limiting functions span $S$ times polynomials of degree at
most two in
\[
 \mu=\frac{\rho^2}{1+\rho^2}.
\]
Their value and first derivative at $\alpha$, together with their value
at the distinct point $\beta$, form a nonsingular Hermite matrix; the
change from $\theta$ to $\mu$ has nonzero derivative. Undoing the
column operations shows that the normalized three-column determinant
has magnitude comparable to $L^{-3}$. The normalized scalar denominator
tends to $S(\rho(\alpha))\ne0$.

In particular the confluent denominator $\det K_{S_1}$ does not vanish
at large even sizes. The preceding row and column identities then also
give $Q_o(1,\ldots,1)\ne0$, so the rational crossing formula may indeed
be specialized there. Undoing the row and column normalizations and
canceling the scalar denominator yields
\[
 {\cal B}_n(1,\ldots,1)
 \asymp L^{-3}A_L(2)A_L(3)f_L(\alpha)f_L(\beta)
 \asymp n^{-1/4}.
\]
Here positivity comes from the convergent strip comparison,
$\rho(\alpha)\rho(\beta)=1$, and the remaining power is
$-3+h_2+h_3=-3+1+7/4=-1/4$. Finally the positive column inequalities
and monotonicity in honeycomb width extend the two-sided bound from
even square sizes to all integer honeycomb widths:
${\cal B}_n\asymp n^{-1/4}$. This proves Proposition~\ref{L:prop:strip-mass}.
